\section{Introduction}
An embedded spiking controller must produce an effective action before contact. A neural-rate criterion can therefore accept an episode that fails its task, while altered network activity can accompany successful behavior. Assessment requires a defined plant, a causal neural-to-actuator interface and a declared observation horizon.

Closed-loop comparisons also couple faults to future sensory input. Equal seed labels alone need not align random inputs when stopping or event consumption changes. Early termination changes the interval over which spikes are accumulated. Stopped-episode, common-duration and fixed-horizon counts answer different resource questions. Their relationship to task failure should be measured on paired records after verifying healthy capability in every direction.

We couple a fly-connectome spiking substrate to a synthetic pursuit-avoidance task. Its assumed receptive fields and body coupling define an engineering stress model, not biological escape. Separate development and validation establish two bilateral laws; new seed clusters cross law, direction and nested deletion dose at a fixed one-second endpoint.

Our engineering contribution is an executable paired protocol with source-bound registrations, shared input variates and masks, exact actuation and independent input/geometry replay. The prospective experiment demonstrates controller-dependent fault loss despite identical healthy success counts, and quantifies errors when neural criteria are used as task-failure warnings. Healthy safety reserves and executed actions still differ; the comparison concerns the whole system. Standard statistical and kinematic components support these bounded findings, not new algorithms. All failed development attempts and test cells are retained.
